% ============================================
% PERSONAL INFORMATION - SINGLE SOURCE OF TRUTH
% ============================================

% Basic Info
\newcommand{\myname}{Harshit Kandpal}
\newcommand{\myemail}{kandpalhar@gmail.com}
\newcommand{\mygithub}{HarK-github}
\newcommand{\mylinkedin}{harshit-k-a746a1310}
\newcommand{\myphone}{+91 9321763273}
\newcommand{\mylocation}{Durg, India (IST/UTC+5:30)}
\newcommand{\mydegree}{B.Tech Computer Science, IIT Bhilai (2024–2028)}

% Project Info
\newcommand{\projecttitle}{Improve usability, maintainability, and contributor experience of Hyperledger Cacti}
\newcommand{\program}{LF Decentralized Trust Mentorship 2026}
\newcommand{\mentors}{André Augusto, Carlos Amaro, Rafael Belchior, and Venkatraman Ramakrishna}
\newcommand{\committee}{LF Decentralized Trust Mentorship Program Committee}

% Links
\newcommand{\githuburl}{https://github.com/\mygithub}
\newcommand{\linkedinurl}{https://www.linkedin.com/in/\mylinkedin}

% ============================================
% COMMANDS FOR DISPLAYING INFO (consistent styling)
% ============================================

% Header with icons (small version for sincerity section)
\newcommand{\personalinfoinline}{%
    {\color{secondary}\small \texttt{\myemail}} \quad
    {\color{primary}\small\textbullet} \quad
    {\color{primary}\small \href{\githuburl}{\color{primary}\mygithub}} \quad
    {\color{secondary}\small\textbullet} \quad
    {\color{secondary}\small \href{\linkedinurl}{\color{secondary}\mylinkedin}} \quad
    {\color{primary}\small\textbullet} \quad
    {\color{primary}\small \texttt{\myphone}} \quad
    {\color{secondary}\small\textbullet} \quad
    {\color{secondary}\small \texttt{\mylocation}}
}

% Full header with big name (for cover letter top)
\newcommand{\coverletterheader}{%
    \begin{center}
        {\color{primary}\fontsize{28}{0}\selectfont \textbf{\myname}}
        \vspace{0.2cm}

        \personalinfoinline
    \end{center}
}


% Reduce top margin
\vspace*{-1.5cm}

% Add full-width shaded background banner ONLY on this page
\AddToShipoutPictureBG*{%
    \color{lightgrayblue}
    \AtPageUpperLeft{\rule[-1.7in]{\paperwidth}{1.7in}}
}

% Header (centered, big name, links below)
\coverletterheader

\vspace{0.8cm}

% Cover letter title (BELOW background, bigger text)
\begin{center}
    {\color{primary}\LARGE \textbf{Cover Letter: \projecttitle}}
\end{center}

\vspace{0.5cm}

% Salutation (BELOW background, bigger text)
{\color{primary}\large \textbf{To:}} \committee \\
{\color{secondary}\large \textbf{Mentors:}} \mentors \\
{\color{primary}\large \textbf{Date:}} May 2026

\vspace{0.8cm}

Hyperledger Cacti represents a pivotal advancement in blockchain interoperability, successfully merging the legacies of Hyperledger Cactus and Weaver. However, as the project matures, the "fragmentation debt" from these two distinct architectures remains a significant hurdle for usability and maintenance. The current developer experience is often hampered by long CI/CD cycles (averaging 45 minutes) and a documentation landscape that remains split between "Node Server" and "Relay" concepts. The challenge lies in harmonizing these legacies into a unified, high-performance ecosystem that is easy to contribute to and simple to deploy.

Improving the contributor experience in Cacti is not merely about moving files; it requires a strategic overhaul of the project's infrastructure and knowledge base. This project aims to address these core inefficiencies through a two-pronged approach: optimizing the CI/CD pipeline and restructuring the documentation into a clear, tier-based architecture. By implementing advanced caching strategies (such as \texttt{node\_modules} and \texttt{.tsbuildinfo} persistence) and parallelizing execution, we can reduce CI runtimes by over 60\%, drastically increasing the velocity of the development cycle. Simultaneously, unifying the documentation into structured Level 1 (Library), Level 2 (API/REST), and Level 3 (Framework) tiers will provide a clear path for newcomers, regardless of their background in Cactus or Weaver.

My technical background is uniquely suited to the "Cleanup Initiative." Over the past few years, I have developed deep expertise in TypeScript monorepos, Docker orchestration, and DevOps optimization. I have already invested significant effort into auditing Cacti’s current state, producing detailed reports on CI/CD bottlenecks and documentation fragmentation. These audits (available in the \texttt{researches/} folder) provide a data-driven foundation for the work proposed here. My experience in containerizing multi-ledger systems and my familiarity with the Hyperledger ecosystem ensure that I can hit the ground running and deliver immediate, measurable improvements to the Cacti codebase.

I am particularly excited to work with André Augusto, Carlos Amaro, Rafael Belchior, and Venkatraman Ramakrishna. Their leadership in the Cacti and Weaver communities is inspiring, and I am eager to contribute to their vision of a unified interoperability framework. My goal is to deliver more than just a set of optimizations; I want to help establish a new standard for maintainability and contributor experience within the Hyperledger ecosystem. I am confident that my technical skills, proactive research, and dedication to open-source excellence make me an ideal candidate for this mentorship.

% ==================== ANSWERS TO QUESTIONS ====================
\section*{How I Found Out About the Mentorship Program}

I discovered the OpenSSF LFX Mentorship Program directly through the LFDT website, where I actively search for structured open-source opportunities. The program's focus on impactful projects within supply chain security, like the Hyperledger Fabric-X and Fablo integration, immediately captured my attention. I find these programs crucial for gaining hands-on experience in real-world infrastructure rather than isolated academic exercises.

% Divider line after question 1
\begin{center}
    \textcolor{secondary}{\rule{0.7\textwidth}{0.3pt}}
\end{center}

\section*{Why I Am Interested in This Program}

My interest stems from a long-standing engagement with open-source and decentralized technologies, seeking to apply my skills to impactful hands-on development. The Fablo project's role in lowering the barrier to entry for Hyperledger Fabric-X deeply resonates with me, highlighting its community importance. This ambitious integration offers a significant learning curve, promising substantial growth in my understanding of complex systems and the satisfaction of contributing to a tool that benefits many.

% Divider line after question 2
\begin{center}
    \textcolor{secondary}{\rule{0.7\textwidth}{0.3pt}}
\end{center}

\section*{Experience and Knowledge/Skills Applicable to This Program}

My two years of open-source contributions, marked by over 800 GitHub contributions, have provided extensive experience in distributed systems, Docker-based tooling, TypeScript/Node.js, and CI/CD workflows, all directly relevant to this project. I have leveraged Docker and GitHub Actions for containerized applications and multi-service orchestration, crucial for managing Fabric-X networks via Docker Compose and Fablo. Specifically, my prior study of Fablo’s architecture and the Fabric-X token sample, including its xdev topology, core.yaml, routing tables, and fxconfig lifecycle, culminated in a functional MVP implementing `global.platform: "fabricx"` generator branching and network lifecycle commands. This focused work positions me exceptionally well to drive the Fabric-X integration into Fablo.

% Divider line after question 3
\begin{center}
    \textcolor{secondary}{\rule{0.7\textwidth}{0.3pt}}
\end{center}

\section*{What I Hope to Get Out of This Mentorship Experience}

Through this mentorship, I aim to significantly enhance my proficiency with complex codebase architectures and deepen my understanding of both Fablo and Fabric-X. I particularly look forward to gaining advanced confidence in Docker-based environments and microservice-oriented systems. Learning directly from the mentors' real-world experience and guidance is paramount for my development as an effective open-source contributor, enabling me to make lasting, meaningful contributions to the community and the project's future evolution.

% Divider line after question 4
\begin{center}
    \textcolor{secondary}{\rule{0.7\textwidth}{0.3pt}}
\end{center}

\vspace{0.5cm}

% Sincerity section (left-aligned, bigger name)
\begin{flushleft}
    \textbf{\color{primary}Sincerely,} \\
    \vspace{0.2cm}
    {\color{primary}\Large \textbf{Harshit Kandpal}} \\
    \vspace{0.1cm}
    {\color{secondary}B.Tech Computer Science, IIT Bhilai (2024–2028)}
\end{flushleft}